\documentclass{article}
\usepackage{amsmath,amssymb,geometry}
\usepackage[UTF8]{ctex}
\begin{document}

\section*{P183：以一致的运输单位来完整构建入境物流模型}
这个经过独立修复的版本保留了原稿中的主要内容、目录、交付模式以及容量系列等要素。同时，它修正了提取出的索引以及过于复杂的价格/重量表示方式，采用了下面列出的明确索引和含义。

\subsection*{输入与变量}
供应商为$i\in L$，日期为$j\in D=\{1,\ldots,N\}$，其中正数偶数的数量为$N$；模式为$M=\{\mathrm{FTL},\mathrm{CES},\mathrm{LTL}\}$，频率为$o\in O$；区域为$z\in Z$，存储资源为$h\in H$，LTL重量等级为$b\in Q$，CES等级为$k\in K$。$d_i>0$表示交货量；$q$和$s$则是该交货量的分数部分。$SS_i$是安全库存的比例。$f_i$是每单位体积的货物重量。$e_i$是空容器体积系数。$\omega_i$则是每单位空容器体积的重量。所有关税系数$c^L_{ib},c^C_k$都与重量等级的限制$(L^L_{ib},U^L_{ib})$和$(L^C_k,U^C_k)$分开计算。
零CES重量等级的关税为零。实施过程中的JSON名称分别为fwq、fSLC、omega\_ec、Bp、BpCES、WeightLo/WeightHi以及CESLo/CESHi。

二进制变量为$v_i^m,p_{ij}^m,\gamma_i,\tau_{ij},\beta_{io},
\alpha_{bij},\alpha^{ec}_{bij},\delta_{kij}$。  
$0\le q_{ij}^m\le1$、$s_{ij}\ge0$（包括$s_{i0}$）以及$w_{bij},w^{ec}_{bij},w^C_{kij}\ge0$。  
$n_{ij},n^{ec}_{ij},n^L_{jz}$是非负整数。  
以下所有求和运算都适用于所声明的索引集；由半周期索引的约束条件则使用$j=1,\ldots,N/2$来表示。

\subsection*{目标}
\begin{align}
\min\quad&
\sum_{i,j}C_i^D(n_{ij}+n^{ec}_{ij})
+\sum_{i,j,k}c_k^C\delta_{kij}
+\sum_{i,j,b}c^L_{ib}(w_{bij}+w^{ec}_{bij})\\
&+A\sum_{i,j,m}p_{ij}^m
+N\sum_{i,o}(C_i^{dR}u_{io}^R+C_i^{dI}u_{io}^I)\beta_{io}.
\end{align}
如论文第6页方程(1)所示，这五个成本组成部分被相加在一起。
\subsection*{需求、库存和模式选择}
\begin{align}
\sum_{j,m}q_{ij}^m&=1,&
s_{i0}+\sum_{\ell=1}^j\sum_mq_{i\ell}^m&\ge j/N,\\
s_{ij}&=s_{i,j-1}+\sum_mq_{ij}^m-1/N+SS_i\mathbf1_{\{j=1\}},\\
\sum_m v_i^m&=1,&
\sum_m p_{ij}^m&\le1,&
q_{ij}^m&\le p_{ij}^m\le v_i^m,\\
q_{ij}^m&\ge(Q^{min}/d_i)p_{ij}^m\quad(m=\mathrm{FTL},\mathrm{CES}),&
q_{ij}^{LTL}&\ge p_{ij}^{LTL}/N.
\end{align}
初始库存比例仍然是一个决策变量，就像在所选模型中一样。

\subsection*{循环模式与频率}
\begin{align}
-\gamma_i&\le p_{ij}^m-p_{i,j+N/2}^m\le\gamma_i,\\
\sum_{j,m}p_{ij}^m&\le1+N(1-\gamma_i),\\
\sum_{j=1}^{N/2}q_{ij}^{LTL}&\le
 \sum_{j=1}^{N/2}q_{i,j+N/2}^{LTL}+\gamma_i,\\
\sum_mq_{ij}^m&\ge\tau_{ij}/N,&
\sum_j\tau_{ij}&=\sum_o o\beta_{io},&
\sum_o\beta_{io}&=1,\\
\sum_{j=1}^{N/2}\sum_m p_{ij}^m&\ge1-\gamma_i,&
\sum_{j=N/2+1}^{N}\sum_m p_{ij}^m&\ge1-\gamma_i,&
\sum_jp_{ij}^m&\ge v_i^m.
\end{align}
这些不等式定义了$\tau$与交付指标之间的单向关系。

\subsection*{入口、储存空间以及车辆的容量}
\begin{align}
\sum_i n_{ij}+\sum_z n^L_{jz}&\le Cap_{GI},&
\sum_i p_{ij}^{CES}&\le Cap_{WK},\\
\sum_i fqp_{hi}d_i s_{ij}&\le Cap_h,&
d_iq_{ij}^{FTL}&\le Cap_L n_{ij},\\
f_i d_iq_{ij}^{FTL}&\le Cap_{WL}n_{ij},&
e_id_iq_{ij}^{FTL}&\le Cap_Ln^{ec}_{ij},\\
e_i\omega_i d_iq_{ij}^{FTL}&\le Cap_{WL}n^{ec}_{ij},&
f_id_iq_{ij}^{CES}&\le Cap_K,\\
f_id_iq_{ij}^{LTL}&\ge\omega^{LTL}p_{ij}^{LTL},\\
\sum_i r_{iz}d_iq_{ij}^{LTL}&\le Cap_Ln^L_{jz},&
\sum_i r_{iz}f_id_iq_{ij}^{LTL}&\le Cap_{WL}n^L_{jz}.
\end{align}
FTL空容器重量约束现在包含$d_i$，与LTL空容器重量相同。这是唯一的第二轮求解器代码修正。

\subsection*{货运及空运回程的重量等级}
\begin{align}
\sum_b\alpha_{bij}&=1,&
L^L_{ib}\alpha_{bij}&\le w_{bij}\le U^L_{ib}\alpha_{bij},&
\sum_bw_{bij}&=f_id_iq_{ij}^{LTL},\\
\sum_b\alpha^{ec}_{bij}&=1,&
L^L_{ib}\alpha^{ec}_{bij}&\le w^{ec}_{bij}\le U^L_{ib}\alpha^{ec}_{bij},&
\sum_bw^{ec}_{bij}&=e_i\omega_i d_iq_{ij}^{LTL},\\
\sum_k\delta_{kij}&=1,&
L^C_k\delta_{kij}&\le w^C_{kij}\le U^C_k\delta_{kij},&
\sum_kw^C_{kij}&=f_i d_iq_{ij}^{CES}.
\end{align}

\subsection*{提供并修复证据}
论文第6-8页以及第15-16页中，方程(6)、(17)、(18)、(38)至(49)分别区分了 fractional shipment quantity与物理容量单位。第一次修复将价格比率与重量限制分开处理，恢复了$b$在$w$上的缺失索引，并且恢复了$Q^{min}/d_i$、$SS_i$以及存储因子$d_i$。第二次修复则修正了剩余的FTL空重量因子。对于当前所有的十个实例，$\omega_i Cap_L\le Cap_{WL}$和$d_i\ge1$都参与其中；因此，无论是之前的弱重量不等式，还是经过修正的不等式，都可以通过空体积容量来理解。这就解释了为什么在这些输入条件下，之前的值能够被保留下来；不过，修正后的代码仍然需要对所有十个实例进行重新运行。

\end{document}
